
\section{Use Cases}
\label{sec:usecases}

In this section, we discuss two representative use cases:
(1) private explanations for query refinement, and
(2) private explanations for consistent query answering (CQA).

\subsection{Query Refinement under Diversity Constraints}

One of the primary motivations for this work is the growing need to enforce fairness,
representation, and diversity constraints in data-driven systems.
Recent works demonstrated that ranking, filtering, and query-processing systems
may unintentionally underrepresent protected demographic groups
or amplify existing societal biases \cite{barocas2023,friedler2016,stoyanovich2018}.

To address these concerns, several works proposed mechanisms for enforcing fairness
and diversity constraints directly during query processing and ranking generation
\cite{stoyanovich2018,erica2023}.
In particular, the \textsc{ERICA} framework \cite{erica2023}
studied the problem of refining SQL queries in order to satisfy
cardinality constraints over demographic groups while minimizing
semantic deviation from the original query intent.

Consider the following hiring query:

\begin{verbatim}
Q:
SELECT *
FROM Candidates
WHERE Experience >= 5
AND SalaryExpectation <= 20000
AND Degree = 'CS'
\end{verbatim}

Suppose the result of this query contains very few women candidates,
violating a diversity constraint requiring that at least $30\%$
of the returned candidates belong to the group
\texttt{Gender = Female}.

A query refinement algorithm such as \textsc{ERICA} may generate
the following refined query:

\begin{verbatim}
Q':
SELECT *
FROM Candidates
WHERE Experience >= 4
AND SalaryExpectation <= 24000
AND Degree IN ('CS', 'Software Engineering')
\end{verbatim}

The refined query satisfies the diversity constraints,
but it also changes the semantic meaning of the result set.
Existing query-refinement systems mainly focus on generating the refined query itself,
with relatively limited support for explaining the semantic impact
of the refinement process \cite{mishra2009,erica2023}.

Our framework naturally complements these systems by generating
privacy-preserving explanations for the differences between $Q$
and $Q'$.
For example, our algorithm may identify that the refinement
primarily changes the distributions of:
\begin{itemize}
    \item education level,
    \item years of experience,
    \item expected salary,
    \item and employment history.
\end{itemize}

The algorithm may then release noisy histograms showing that:
\begin{itemize}
    \item the refined query includes substantially more candidates
    with $4$--$5$ years of experience,
    \item candidates from additional engineering disciplines now appear,
    \item and the salary distribution shifted toward higher salary expectations.
\end{itemize}

Importantly, these explanations are generated under differential privacy,
ensuring that no individual candidate can be inferred from the explanation outputs.

Such explanations provide several practical benefits:
\begin{enumerate}
    \item helping users understand the semantic consequences of the refinement,
    \item enabling analysts to evaluate whether the refined query still aligns
    with the original intent,
    \item identifying unintended side effects,
    \item and improving transparency in fairness-aware systems.
\end{enumerate}

More generally, our framework can serve as an interpretability layer
on top of existing query-refinement mechanisms.
Because the framework operates over arbitrary query pairs $(Q,Q')$,
it can be integrated with many refinement strategies,
including interactive query refinement systems \cite{mishra2009},
fairness-aware ranking systems \cite{stoyanovich2018},
and provenance-based bias mitigation frameworks \cite{moskovitch2022}.

\subsection{Consistent Query Answering}

Another important application domain for our framework is
\emph{consistent query answering} (CQA).
The goal of CQA is to provide meaningful answers to queries
over databases that may violate integrity constraints
\cite{bertossi2011,bertossi2006}.
Instead of querying the inconsistent database directly,
CQA systems construct repairs or consistent versions of the data
and evaluate the query with respect to those repaired instances.

The CQA literature studied many forms of integrity constraints,
including functional dependencies, denial constraints,
key constraints, and inclusion dependencies
\cite{bertossi2011,arenas1999}.
Different repair strategies may lead to substantially different query answers,
raising important interpretability and transparency questions.

Consider the following example involving employee records:

\begin{verbatim}
Employees(EmpID, Department, Salary)
\end{verbatim}

Suppose the database violates a key constraint because multiple tuples
exist with the same \texttt{EmpID} but different salary values.
A repair algorithm may remove conflicting tuples or choose a canonical tuple.

Now consider the following query:

\begin{verbatim}
Q:
SELECT Department, AVG(Salary)
FROM Employees
GROUP BY Department
\end{verbatim}

Let $Q'$ denote the same query evaluated over the repaired database instance.
The repaired answer may differ significantly from the original answer,
especially if the inconsistencies are concentrated in specific departments.

Our framework can provide private explanations for the differences
between $Q$ and $Q'$.
For example, the algorithm may identify that:
\begin{itemize}
    \item the Finance department experienced the largest salary changes,
    \item duplicate tuples were concentrated among senior employees,
    \item and the repaired database substantially altered the salary distribution
    of employees with more than ten years of experience.
\end{itemize}

The explanation mechanism may release noisy histograms comparing
the salary distributions before and after the repair process.
These explanations help analysts understand the semantic effects
of the repair strategy without revealing information about specific employees.

This capability is particularly important because traditional explanation methods
for CQA may inadvertently leak sensitive information.
When repairs are small or localized,
revealing exact changes may expose individual tuples or conflicting records.
By contrast, our framework provides formal differential privacy guarantees
\cite{dwork2014}, ensuring that the explanations themselves
preserve privacy.

Our approach also enables comparison between different repair algorithms.
For example, two repair mechanisms may both satisfy the integrity constraints,
yet one may disproportionately affect records associated with a particular region,
demographic group, or salary range.
Our explanation framework enables analysts to compare these semantic effects
through differentially private summaries.

More broadly, the generalized formulation of our framework
as a mechanism for privately explaining differences between query outputs
makes it applicable to a wide variety of data-management tasks,
including:
\begin{itemize}
    \item fairness-aware query processing,
    \item data cleaning,
    \item provenance analysis,
    \item query debugging,
    \item and consistent query answering.
\end{itemize}

This generalization extends the original query-refinement motivation
into a broader framework for interpretable and privacy-preserving
data analytics.
